% Options for packages loaded elsewhere
% Options for packages loaded elsewhere
\PassOptionsToPackage{unicode}{hyperref}
\PassOptionsToPackage{hyphens}{url}
\PassOptionsToPackage{dvipsnames,svgnames,x11names}{xcolor}
%
\documentclass[
  ngerman,
  10pt,
  twocolumn]{article}
\usepackage{xcolor}
\usepackage[top=2.5cm,bottom=2.5cm,left=2cm,right=2cm]{geometry}
\usepackage{amsmath,amssymb}
\setcounter{secnumdepth}{-\maxdimen} % remove section numbering
\usepackage{iftex}
\ifPDFTeX
  \usepackage[T1]{fontenc}
  \usepackage[utf8]{inputenc}
  \usepackage{textcomp} % provide euro and other symbols
\else % if luatex or xetex
  \usepackage{unicode-math} % this also loads fontspec
  \defaultfontfeatures{Scale=MatchLowercase}
  \defaultfontfeatures[\rmfamily]{Ligatures=TeX,Scale=1}
\fi
\usepackage{lmodern}
\ifPDFTeX\else
  % xetex/luatex font selection
\fi
% Use upquote if available, for straight quotes in verbatim environments
\IfFileExists{upquote.sty}{\usepackage{upquote}}{}
\IfFileExists{microtype.sty}{% use microtype if available
  \usepackage[]{microtype}
  \UseMicrotypeSet[protrusion]{basicmath} % disable protrusion for tt fonts
}{}
\usepackage{setspace}
\makeatletter
\@ifundefined{KOMAClassName}{% if non-KOMA class
  \IfFileExists{parskip.sty}{%
    \usepackage{parskip}
  }{% else
    \setlength{\parindent}{0pt}
    \setlength{\parskip}{6pt plus 2pt minus 1pt}}
}{% if KOMA class
  \KOMAoptions{parskip=half}}
\makeatother
% Make \paragraph and \subparagraph free-standing
\makeatletter
\ifx\paragraph\undefined\else
  \let\oldparagraph\paragraph
  \renewcommand{\paragraph}{
    \@ifstar
      \xxxParagraphStar
      \xxxParagraphNoStar
  }
  \newcommand{\xxxParagraphStar}[1]{\oldparagraph*{#1}\mbox{}}
  \newcommand{\xxxParagraphNoStar}[1]{\oldparagraph{#1}\mbox{}}
\fi
\ifx\subparagraph\undefined\else
  \let\oldsubparagraph\subparagraph
  \renewcommand{\subparagraph}{
    \@ifstar
      \xxxSubParagraphStar
      \xxxSubParagraphNoStar
  }
  \newcommand{\xxxSubParagraphStar}[1]{\oldsubparagraph*{#1}\mbox{}}
  \newcommand{\xxxSubParagraphNoStar}[1]{\oldsubparagraph{#1}\mbox{}}
\fi
\makeatother


\usepackage{longtable,booktabs,array}
\usepackage{calc} % for calculating minipage widths
% Correct order of tables after \paragraph or \subparagraph
\usepackage{etoolbox}
\makeatletter
\patchcmd\longtable{\par}{\if@noskipsec\mbox{}\fi\par}{}{}
\makeatother
% Allow footnotes in longtable head/foot
\IfFileExists{footnotehyper.sty}{\usepackage{footnotehyper}}{\usepackage{footnote}}
\makesavenoteenv{longtable}
\usepackage{graphicx}
\makeatletter
\newsavebox\pandoc@box
\newcommand*\pandocbounded[1]{% scales image to fit in text height/width
  \sbox\pandoc@box{#1}%
  \Gscale@div\@tempa{\textheight}{\dimexpr\ht\pandoc@box+\dp\pandoc@box\relax}%
  \Gscale@div\@tempb{\linewidth}{\wd\pandoc@box}%
  \ifdim\@tempb\p@<\@tempa\p@\let\@tempa\@tempb\fi% select the smaller of both
  \ifdim\@tempa\p@<\p@\scalebox{\@tempa}{\usebox\pandoc@box}%
  \else\usebox{\pandoc@box}%
  \fi%
}
% Set default figure placement to htbp
\def\fps@figure{htbp}
\makeatother



\ifLuaTeX
\usepackage[bidi=basic]{babel}
\else
\usepackage[bidi=default]{babel}
\fi
% get rid of language-specific shorthands (see #6817):
\let\LanguageShortHands\languageshorthands
\def\languageshorthands#1{}
\ifLuaTeX
  \usepackage[german]{selnolig} % disable illegal ligatures
\fi


\setlength{\emergencystretch}{3em} % prevent overfull lines

\providecommand{\tightlist}{%
  \setlength{\itemsep}{0pt}\setlength{\parskip}{0pt}}



 
\usepackage[style=authoryear,backend=biber,style=authoryear-comp,uniquename=false,uniquelist=false]{biblatex}
\addbibresource{../references.bib}


\usepackage{csquotes}
\usepackage{booktabs}
\usepackage{array}
\usepackage{multirow}
\usepackage{wrapfig}
\usepackage{float}
\usepackage{colortbl}
\usepackage{pdflscape}
\usepackage{threeparttable}
\usepackage{threeparttablex}
\usepackage[normalem]{ulem}
\usepackage{makecell}
\usepackage{xcolor}
\usepackage{multicol}
\usepackage{fancyhdr}
\usepackage{titlesec}
\usepackage{etoolbox}
\usepackage{tabularx}
\usepackage{afterpage}

% Define WISTA colors (brighter blue)
\definecolor{wistaBlue}{RGB}{30,110,180}
\definecolor{wistaLightBlue}{RGB}{51,153,204}
\definecolor{wistaDarkGray}{RGB}{64,64,64}
\definecolor{wistaLightGray}{RGB}{128,128,128}

% Custom section formatting - number on top with column-width underline
\titleformat{\section}[display]
  {\normalfont\large\bfseries}
  {\color{wistaBlue}\large\bfseries\thesection\vspace{0.1cm}\\\noindent\rule{\columnwidth}{1pt}}
  {0.3cm}
  {\color{wistaDarkGray}\large\bfseries}
  [\vspace{0.1cm}\noindent{\color{wistaBlue}\rule{\columnwidth}{0.5pt}}\vspace{0.3cm}]
\titlespacing*{\section}{0pt}{2ex plus 1ex minus .2ex}{1ex plus .2ex}

% Subsection formatting - number and title with underline
\titleformat{\subsection}
  {\color{wistaDarkGray}\normalsize\bfseries}
  {\color{wistaBlue}\normalsize\bfseries\thesubsection}
  {0.5em}
  {\color{wistaDarkGray}\normalsize\bfseries}
  [\vspace{0.05cm}\noindent{\color{wistaBlue}\rule{0.5\columnwidth}{0.4pt}}\vspace{0.1cm}]
  
% Subsubsection formatting
\titleformat{\subsubsection}
  {\color{wistaDarkGray}\small\bfseries}
  {\color{wistaBlue}\small\bfseries\thesubsubsection}
  {0.3em}
  {\color{wistaDarkGray}\small\bfseries}

% Header and footer setup
\pagestyle{fancy}
\fancyhf{} % Clear all header and footer fields
\fancyhead[LE]{\small\textit{Max Mustermann, Dr. Anna Schmidt}}
\fancyhead[RO]{\small\textit{Effizienzanalyse statistischer Verfahren}}
\fancyfoot[LE]{\small\thepage}
\fancyfoot[RE]{\small Statistisches Bundesamt | WISTA | {\color{wistaBlue}6} | 2025}
\fancyfoot[LO]{\small Statistisches Bundesamt | WISTA | {\color{wistaBlue}6} | 2025}
\fancyfoot[RO]{\small\thepage}
\renewcommand{\headrulewidth}{0pt}
\renewcommand{\footrulewidth}{0pt}

% Make all figures and tables span both columns (full width)
% Redefine figure environment to use figure*
\let\origfigure\figure
\let\endorigfigure\endfigure
\renewenvironment{figure}[1][htbp]{%
  \origfigure*[#1]%
}{%
  \endorigfigure*%
}

% Redefine table environment to use table*
\let\origtable\table
\let\endorigtable\endtable
\renewenvironment{table}[1][htbp]{%
  \origtable*[#1]%
}{%
  \endorigtable*%
}
\makeatletter
\@ifpackageloaded{caption}{}{\usepackage{caption}}
\AtBeginDocument{%
\ifdefined\contentsname
  \renewcommand*\contentsname{Inhaltsverzeichnis}
\else
  \newcommand\contentsname{Inhaltsverzeichnis}
\fi
\ifdefined\listfigurename
  \renewcommand*\listfigurename{Abbildungsverzeichnis}
\else
  \newcommand\listfigurename{Abbildungsverzeichnis}
\fi
\ifdefined\listtablename
  \renewcommand*\listtablename{Tabellenverzeichnis}
\else
  \newcommand\listtablename{Tabellenverzeichnis}
\fi
\ifdefined\figurename
  \renewcommand*\figurename{Abbildung}
\else
  \newcommand\figurename{Abbildung}
\fi
\ifdefined\tablename
  \renewcommand*\tablename{Tabelle}
\else
  \newcommand\tablename{Tabelle}
\fi
}
\@ifpackageloaded{float}{}{\usepackage{float}}
\floatstyle{ruled}
\@ifundefined{c@chapter}{\newfloat{codelisting}{h}{lop}}{\newfloat{codelisting}{h}{lop}[chapter]}
\floatname{codelisting}{Listing}
\newcommand*\listoflistings{\listof{codelisting}{Listingverzeichnis}}
\makeatother
\makeatletter
\makeatother
\makeatletter
\@ifpackageloaded{caption}{}{\usepackage{caption}}
\@ifpackageloaded{subcaption}{}{\usepackage{subcaption}}
\makeatother
\usepackage{bookmark}
\IfFileExists{xurl.sty}{\usepackage{xurl}}{} % add URL line breaks if available
\urlstyle{same}
\hypersetup{
  pdftitle={Effizienzanalyse statistischer Verfahren in der amtlichen Statistik},
  pdfauthor={Max Mustermann; Dr.~Anna Schmidt},
  pdflang={de},
  colorlinks=true,
  linkcolor={blue},
  filecolor={Maroon},
  citecolor={Blue},
  urlcolor={Blue},
  pdfcreator={LaTeX via pandoc}}


\title{Effizienzanalyse statistischer Verfahren in der amtlichen
Statistik}
\usepackage{etoolbox}
\makeatletter
\providecommand{\subtitle}[1]{% add subtitle to \maketitle
  \apptocmd{\@title}{\par {\large #1 \par}}{}{}
}
\makeatother
\subtitle{Eine empirische Untersuchung am Beispiel der
Datenverarbeitung}
\author{Max Mustermann \and Dr.~Anna Schmidt}
\date{2025-11-22}
\begin{document}
\maketitle


\setstretch{1.1}
\subsection{Infrastruktur des FDZ-Bund}\label{sec-infrastruktur}

\subsubsection{Die Forschungsdatenzentren der amtlichen
Statistik}\label{die-forschungsdatenzentren-der-amtlichen-statistik}

Die Forschungsdatenzentren (FDZ) der amtlichen Statistik bestehen aus
zwei organisatorisch voneinander getrennten Einrichtungen: dem
\textbf{Forschungsdatenzentrum des Statistischen Bundesamts} und dem
\textbf{Forschungsdatenzentrum der Statistischen Ämter der Länder}.
Beide FDZ arbeiten eng zusammen und bieten ein abgestimmtes Daten- und
Dienstleistungsangebot für die wissenschaftliche Nutzung von Mikrodaten
der amtlichen Statistiken.

\paragraph{Entstehung und Entwicklung}\label{entstehung-und-entwicklung}

Die Gründung der FDZ geht auf ein Gutachten der „Kommission zur
Verbesserung der informationellen Infrastruktur zwischen Wissenschaft
und Statistik'' (KVI) von 1999 zurück. Mit finanzieller Förderung des
Bundesministeriums für Bildung und Forschung (BMBF) wurde im Herbst 2001
zunächst das FDZ des Statistischen Bundesamts gegründet. Das FDZ der
Statistischen Ämter der Länder folgte im April 2002.

\paragraph{Aufgaben und Ziele}\label{aufgaben-und-ziele}

Die FDZ verfolgen das Ziel, den Zugang zu Mikrodaten der amtlichen
Statistik für die Wissenschaft nicht nur zu ermöglichen, sondern
fortlaufend zu verbessern und an die sich ändernden Bedarfe anzupassen:

\begin{itemize}
\tightlist
\item
  \textbf{Datenzugang:} Bereitstellung hochwertiger Mikrodaten für die
  wissenschaftliche Forschung
\item
  \textbf{Datenschutz:} Gewährleistung des Schutzes personenbezogener
  und unternehmensbezogener Daten nach § 16 BStatG
\item
  \textbf{Beratung:} Unterstützung der Forschenden bei der Datennutzung
  und methodischen Fragen
\item
  \textbf{Infrastruktur:} Bereitstellung einer sicheren und
  leistungsfähigen Auswertungsumgebung
\item
  \textbf{Regionalisierung:} Flächendeckende Standorte in ganz
  Deutschland für kurze Anfahrtswege
\end{itemize}

\subsubsection{Technische Infrastruktur}\label{technische-infrastruktur}

\paragraph{Hardware-Ausstattung}\label{hardware-ausstattung}

Die technische Infrastruktur des FDZ-Bund umfasst verschiedene
Systemkomponenten:

\begin{table}

\caption{\label{tbl-fdz-infrastructure}Technische Ausstattung des
FDZ-Bund (beispielhafte Konfiguration)}

\centering{

\begin{tabular}{llc}
\toprule
Komponente & Spezifikation & Anzahl\\
\midrule
Rechenknoten (Standard) & 2x Intel Xeon Gold 6248R (48 Kerne), 384 GB RAM & 12\\
Rechenknoten (High-Memory) & 2x Intel Xeon Platinum 8280 (56 Kerne), 1,5 TB RAM & 4\\
Speichersystem & NAS-System mit 500 TB nutzbarem Speicher & 1\\
Netzwerk & 10 Gbit/s Ethernet, InfiniBand für parallele Operationen & -\\
Backup-System & Tägliches Backup, räumlich getrennte Archivierung & -\\
\bottomrule
\end{tabular}

}

\end{table}%

\paragraph{Software-Umgebung}\label{software-umgebung}

Den Forschenden stehen verschiedene Software-Tools zur Verfügung:

\begin{itemize}
\tightlist
\item
  \textbf{Statistische Software:} R, Python, Stata, SAS, SPSS
\item
  \textbf{Datenbanksysteme:} PostgreSQL, DuckDB
\item
  \textbf{Entwicklungsumgebungen:} RStudio Server, JupyterLab
\item
  \textbf{Versionskontrolle:} Git (lokal, ohne externe Verbindung)
\end{itemize}

\subsubsection{Infrastrukturkonzept}\label{infrastrukturkonzept}

\paragraph{Fachlich zentralisierte
Datenhaltung}\label{fachlich-zentralisierte-datenhaltung}

In Deutschland wird der überwiegende Teil der amtlichen
Statistikproduktion dezentral in den Statistischen Ämtern der Länder
durchgeführt (über 90\% aller amtlichen Statistiken). Da sich
wissenschaftliche Analysen in der Regel jedoch auf mehrere Bundesländer
oder das gesamte Bundesgebiet beziehen, haben die FDZ eine
\textbf{fachlich zentralisierte Datenhaltung} eingerichtet. Diese
ermöglicht es, die Daten des gesamten Bundesgebiets länderübergreifend
an jedem Standort bereitzustellen.

\paragraph{Regionalisierte
Infrastruktur}\label{regionalisierte-infrastruktur}

Die FDZ verfolgen das Ziel, immer in der Nähe der Wissenschaft zu sein.
Dies wird durch über ganz Deutschland verteilte Standorte erreicht.
Standorte befinden sich u.a. in:

\begin{itemize}
\tightlist
\item
  Berlin, Bonn, Wiesbaden (Bund)
\item
  Bad Ems, Bremen, Dresden, Düsseldorf, Erfurt, Frankfurt, Hamburg,
  Hannover, Kiel, München, Stuttgart u.v.m. (Länder)
\end{itemize}

\subsubsection{Zugangsformen und
Nutzungsmodelle}\label{zugangsformen-und-nutzungsmodelle}

Die FDZ bieten verschiedene Zugangswege an, über die unterschiedlich
anonymisierte Datenprodukte bereitgestellt werden:

\paragraph{Off-Site-Zugang}\label{off-site-zugang}

\begin{itemize}
\tightlist
\item
  \textbf{Public Use Files (PUF):} Stark anonymisierte Daten für freie
  Nutzung
\item
  \textbf{Scientific Use Files (SUF):} Faktisch anonymisierte Daten für
  wissenschaftliche Zwecke
\item
  \textbf{Remote Scientific Use Files (Remote-SUF):} Neuer Zugangsweg
  mit höherer Detailtiefe
\end{itemize}

\paragraph{On-Site-Zugang}\label{on-site-zugang}

\begin{itemize}
\tightlist
\item
  \textbf{Gastwissenschaftsarbeitsplatz (GWAP):} Zugang zu formal
  anonymisierten Daten vor Ort

  \begin{itemize}
  \tightlist
  \item
    Vorteile: Höchste Datendetailtiefe, direkte Beratung
  \item
    Anforderungen: Terminvereinbarung, Zutrittskontrolle,
    Output-Kontrolle
  \end{itemize}
\item
  \textbf{Kontrollierte Datenfernverarbeitung (KDFV):} Remote-Ausführung
  von Analyseskripten

  \begin{itemize}
  \tightlist
  \item
    Vorteile: Keine Anreise, zeitlich flexibel
  \item
    Einschränkungen: Keine interaktive Arbeit, Ergebnisfreigabe
    erforderlich
  \end{itemize}
\end{itemize}

\subsubsection{Datenschutz und
Sicherheit}\label{datenschutz-und-sicherheit}

\paragraph{Mehrstufiges
Sicherheitskonzept}\label{mehrstufiges-sicherheitskonzept}

Das FDZ-Bund implementiert ein umfassendes Sicherheitskonzept:

\begin{enumerate}
\def\labelenumi{\arabic{enumi}.}
\tightlist
\item
  \textbf{Physische Sicherheit:} Zutrittskontrolle, gesicherte
  Räumlichkeiten
\item
  \textbf{Netzwerksicherheit:} Isolierte Netzwerke, Firewalls, Intrusion
  Detection
\item
  \textbf{Zugriffsrechte:} Rollenbasierte Zugriffskontrolle,
  Authentifizierung
\item
  \textbf{Output-Kontrolle:} Prüfung aller Ergebnisse vor Freigabe
\item
  \textbf{Audit-Logging:} Protokollierung aller Zugriffe und Operationen
\end{enumerate}

\paragraph{Output-Kontrolle}\label{output-kontrolle}

Bevor Forschungsergebnisse das FDZ verlassen dürfen, durchlaufen sie
eine mehrstufige Kontrolle:

\begin{itemize}
\tightlist
\item
  \textbf{Automatische Prüfung:} Screening auf kritische Werte (z.B.
  Fallzahlen \textless{} 5)
\item
  \textbf{Manuelle Prüfung:} Sichtung durch geschulte Mitarbeiter
\item
  \textbf{Dokumentation:} Nachvollziehbare Begründung der
  Freigabeentscheidung
\end{itemize}

\subsubsection{Performanz-Optimierung in der
FDZ-Umgebung}\label{performanz-optimierung-in-der-fdz-umgebung}

\paragraph{Herausforderungen}\label{herausforderungen}

Die Arbeit in der geschützten FDZ-Umgebung stellt besondere
Anforderungen:

\begin{itemize}
\tightlist
\item
  \textbf{Begrenzte Ressourcen:} Mehrere Nutzer teilen sich die
  verfügbare Hardware
\item
  \textbf{Keine Cloud-Dienste:} Externe Rechenressourcen können nicht
  genutzt werden
\item
  \textbf{Datengröße:} Große Datensätze müssen effizient verarbeitet
  werden
\item
  \textbf{Reproduzierbarkeit:} Analysen müssen dokumentiert und
  nachvollziehbar sein
\end{itemize}

\paragraph{Optimierungsstrategien}\label{optimierungsstrategien}

Verschiedene Strategien helfen, die Performanz zu optimieren:

\begin{figure}

\centering{

\pandocbounded{\includegraphics[keepaspectratio]{03-infrastruktur_files/figure-pdf/fig-optimization-strategies-1.pdf}}

}

\caption{\label{fig-optimization-strategies}Strategien zur
Performanz-Optimierung im FDZ-Bund}

\end{figure}%

\subsubsection{Akkreditierung und
Qualitätssicherung}\label{akkreditierung-und-qualituxe4tssicherung}

Die Forschungsdatenzentren der Statistischen Ämter des Bundes und der
Länder sind durch den \textbf{Rat für Sozial- und WirtschaftsDaten
(RatSWD) akkreditiert}. Der RatSWD akkreditiert seit 2008
Forschungsdatenzentren, die den Zugang zu qualitativ hochwertigen
sensiblen Daten für die Wissenschaft ermöglichen und dabei die
erforderlichen Kriterien erfüllen.

\subsubsection{Best Practices für
FDZ-Nutzende}\label{best-practices-fuxfcr-fdz-nutzende}

Aus der Erfahrung mit den FDZ lassen sich folgende Empfehlungen für
effizientes Arbeiten ableiten:

\begin{enumerate}
\def\labelenumi{\arabic{enumi}.}
\tightlist
\item
  \textbf{Datenreduktion:} Frühzeitige Filterung auf relevante
  Beobachtungen und Variablen
\item
  \textbf{Inkrementelle Entwicklung:} Erst mit Stichproben arbeiten,
  dann auf Vollsatz skalieren
\item
  \textbf{Effiziente Datenformate:} Nutzung moderner Formate
  (Arrow/Parquet) statt CSV
\item
  \textbf{Dokumentation:} Sorgfältige Dokumentation des Analyseprozesses
  für Output-Kontrolle
\item
  \textbf{Ressourcenplanung:} Bewusste Planung rechenintensiver
  Operationen bei geteilter Infrastruktur
\item
  \textbf{Beratung nutzen:} Frühzeitige Kontaktaufnahme mit
  FDZ-Mitarbeitern bei methodischen Fragen
\end{enumerate}


\printbibliography



\end{document}
